\documentclass{article}

\title{Coupled Mechanical Systems and Linear Algebra}
\author{Jaden Krause}
\date{November 10, 2026}
\usepackage{float}
\usepackage{graphicx}
\usepackage{amsmath}
\makeatletter
\renewcommand{\topfraction}{0.95}       % Allow figures to use up to 95% of the page top
\renewcommand{\textfraction}{0.05}      % Allow text to take as little as 5% of the page 
\renewcommand{\bottomfraction}{0.95}    % Allow figures to use up to 95% of the page bottom
\renewcommand{\floatpagefraction}{0.85}  % Allow a float page to be created even if it's only 20% full
\setcounter{topnumber}{4}               % Allow up to 4 figures at the top of a page
\setcounter{bottomnumber}{4}            % Allow up to 4 figures at the bottom of a page
\setcounter{totalnumber}{8}             % Allow up to 8 total figures per page
\makeatother
\begin{document}
\maketitle

\section{Introduction}
Following our previous analysis of mechanical oscillators, our discussion naturally leads us to further study 
in two areas. The first is coupled mechanical systems. Such a system is defined as having two masses with a 
spring between them and springs on either side of the masses, attached to two supports. The second area of 
interest is linear algebra and the use of matrices in solving such a system. In this paper, we will analyze both 
of these topics, discussing their particular details.

\section{The Nature of Coupled Mechanical Systems}
\begin{figure}
    \centering
    \includegraphics[width=1\textwidth]{FiguresofCoupledSystems/coupled_system_schematic.png}
    \caption{Schematic of a Coupled Mechanical System}
\end{figure}
Coupled mechanical systems can take many different forms, with different support structures and spring arrangements. 
For this discussion, we will be using the structure described in the introduction and illustrated by Figure 1. Furthermore, 
we will be simplifying our model by ignoring the effects of external and internal friction, and air resistance. The 
masses can have different values, and the spring constants can also vary from each other. One characteristic of coupled 
mechanical systems is the normal modes of the systems. A two-mass system, such as the object of our study, has two normal 
modes, moving with either the same direction and magnitude or opposite direction and magnitude. This is illustrated 
in Figure 2. Omega ($\omega$) represents the angular frequency of the vibration.
\begin{figure}
    \centering
    \includegraphics[width=1\textwidth]{FiguresofCoupledSystems/vibrational_modes.png}
    \caption{Schematic of a Coupled Mechanical System}
\end{figure}

\section{Matrices and Coupled Mechanical Systems}
Matrices are useful in that they allow us to organize large groups of information and deal with it concisely and efficiently. 
In a system with motion, as we are dealing with in coupled mechanical systems, they can also be used to represent position and 
velocity. The first matrix of interest to us is the flexibility matrix. Flexibility is measured in $\frac{m}{N}$. Using the 
inverse operation on the flexibility matrix gives us the stiffness matrix for the represented structure. Being the inverse 
of flexibility, stiffness is measured in $\frac{N}{m}$. This operation allows us to deal with the equation:
$$ \textbf{y} = D\textbf{f} $$
Where $\textbf{y}$ represents the displacement and $\textbf{f}$ represents the force applied. Knowing how to use the 
inverse of matrices allows us to solve this equation for both displacement and force, an important capability when dealing 
with these systems. 

\section{Eigenvalues and Eigenvectors}
An important aspect of matrices in coupled mechanical systems regards two related topics: Eigenvalues and Eigenvectors. 
The Eigenvalue ($\lambda = \omega^2$) represents how fast the system oscillates. The Eigenvector represents where the masses 
move relative to each other in the system. The eigenvector is a vector such that when multiplied by associated matrix $A$, it 
results in a change only of magnitude, not direction. The magnitude of the change is the eigenvalue. As shown above, the 
eigenvalue is found using the square of the angular frequency of the system. We use the eigenvalue to find the normal modes, 
which is important because any given oscillation of the system is a combination of its normal modes, even if they look 
random, such as in Figure 3. 
\begin{figure}
    \centering
    \includegraphics[width=1\textwidth]{FiguresofCoupledSystems/displacement_and_time.png}
    \caption{Combination of Normal Modes}
\end{figure}
A physical example of a normal mode is the vibration of a plucked guitar string. The equation to find the eigenvalues is:
$$ \det(D - \lambda I) = 0 $$
Where \textit{D} (the dynamical matrix) is the result of the product $ D = M^{-1}K $. \textit{M} is the mass matrix of 
the system, and \textit{K} is the stiffness matrix. \textit{I} is the identity matrix. Once the eigenvalues are found, 
plug them into the equation $(D - \lambda I)A = 0$ one at a time to find the matching normal mode vector \textit{A}. In 
the section below, we will consider an example.

\section{Mechanical System Example Problem}
$$ M\text{x}'' + K\text{x} = 0 $$
For the symbolic equation listed above, here is the data for our problem:
\begin{align*}
m_1 &= 1 & m_2 &= 1 & (\text{kg}) \\
k_1 &= 2 & k_2 &= 2 & k_3 &= 3 & \left(\frac{\text{N}}{\text{m}}\right)
\end{align*}
We will use this information to find the normal frequencies of the system. This system follows the pattern of the system 
in Figure 1, with two masses connected supports by springs with a spring in between them. First, let us start by identifying 
the mass and stiffness matrices.
\[
M = \begin{bmatrix} 
    m_1 & 0 \\
    0 & m_2
\end{bmatrix} = 
\begin{bmatrix}
    1 & 0 \\
    0 & 1 
\end{bmatrix}
\]
\[
K = \begin{bmatrix}
    k_1 + k_2 & -k_2 \\
    -k_2 & k_2 + k_3
\end{bmatrix} = 
\begin{bmatrix}
    4 & -2 \\
    -2 & 5
\end{bmatrix}
\]
Putting these matrices into the equation $ M\text{x}'' + K\text{x} = 0 $:
\[
\begin{bmatrix}
    1 & 0 \\
    0 & 1
\end{bmatrix}
\begin{bmatrix}
    x_1'' \\
    x_2''
\end{bmatrix}
+ \begin{bmatrix}
    4 & -2 \\
    -2 & 5
\end{bmatrix}
\begin{bmatrix}
    x_1 \\
    x_2
\end{bmatrix}
= \begin{bmatrix}
    0 \\
    0
\end{bmatrix}
\]
Moving on in our analysis, it is necessary to reduce this second-order DE into a linear, first-order problem. Solving for $x''$:
$$ \text{x}'' = -M^{-1}K\text{x} $$
Recall from previously that we use the equation $ \det(D - \lambda I) = 0 $ to find the eigenvalues, which we use to find the 
normal frequencies. To find \textit{D}:
$$ D = M^{-1}K $$
\[
D = \begin{bmatrix}
    1 & 0 \\
    0 & 1
\end{bmatrix}
\begin{bmatrix}
    4 & -2 \\
    -2 & 5 
\end{bmatrix}
= \begin{bmatrix}
    4 & -2 \\
    -2 & 5
\end{bmatrix}
\]
(Since \textit{M} is the identity matrix, its inverse is the same as the original). Now we solve $ \det(D - \lambda I) = 0 $:
\[
\det(D - \lambda I) = \det\begin{bmatrix}
    4-\lambda & -2 \\
    -2 & 5-\lambda 
\end{bmatrix} = 0
\]
$$ (4-\lambda)(5-\lambda)-(-2)(-2)=0 $$
$$ 20-4\lambda-5\lambda+\lambda^2-4=0 $$
$$ \lambda^2-9\lambda+16=0 $$
Using the quadratic formula we obtain:
\begin{align*}
    \lambda_1 &= \frac{9-\sqrt{17}}{2} \\
    \lambda_2 &= \frac{9+\sqrt{17}}{2}
\end{align*}
Using $ \lambda = \omega^2 $, we can find the normal angular frequencies of the system by taking the positive square roots:
\begin{align*}
    \omega_1 &= \sqrt{\frac{9-\sqrt{17}}{2}} \approx 1.562 \text{rad/s} \\
    \omega_2 &= \sqrt{\frac{9+\sqrt{17}}{2}} \approx 2.562 \text{rad/s}
\end{align*}
If we wanted to find the position functions for the system, we would need both the initial values for $x_1'$ and $x_2'$
and the eigenvectors for the system. The equation we use to find the eigenvectors is $ (D-\lambda I)\textbf{v}=0$ where 
\textbf{v} represents the eigenvector. We solve this using each of our eigenvalues:
\[
\begin{bmatrix}
    4-\lambda & -2 \\
    -2 & 5-\lambda
\end{bmatrix}
\begin{bmatrix}
    v_1 \\
    v_2
\end{bmatrix}
= \begin{bmatrix}
    0 \\
    0
\end{bmatrix}
\]
This expression gives us $(4-\lambda)v_1-2v_2=0$, which simplifies to:
$$ 2v_2=(4-\lambda)v_1$$
For $\lambda_1=\frac{9-\sqrt{17}}{2}$,
\begin{align*}
    2v_2&=\left(4-\frac{9-\sqrt{17}}{2}\right)v_1 \\
    2v_2&=\left(\frac{\sqrt{17}-1}{2}\right)v_1    
\end{align*}
Notice that we now have an equation with two variables. To find solutions, we use the characteristic of eigenvectors 
that they can be scaled by any nonzero number. Observe that choosing $v_1=4$ both eliminates the fraction and the 
coefficient of $v_2$.
\begin{align*}
    2v_2&=\left(\frac{\sqrt{17}-1}{2}\right)(4) \\
    v_2&=\sqrt{17}-1\approx 3.12 \\
\end{align*}
Leaving us with:
\[
v_1\approx\begin{bmatrix}
    4 \\
    3.12
\end{bmatrix}
\]
For $\lambda_2=\frac{9+\sqrt{17}}{2}$,
\begin{align*}
    2v_2&=\left(4-\frac{9+\sqrt{17}}{2}\right)v_1 \\
    2v_2&=\left(\frac{-1-\sqrt{17}}{2}\right)v_1
\end{align*}
Notice again that $v_1=4$ clears the denominator and the coefficient of $v_2$.
\begin{align*}
    2v_2&=\left(4-\frac{9+\sqrt{17}}{2}\right)(4) \\
    v_2&=-1-\sqrt{17}\approx -5.12 \\
    v_2&\approx\begin{bmatrix}
        4 \\
        -5.12
    \end{bmatrix}
\end{align*}
Each of these values have tangible correlation to the system under study. The matrix for $v_1$ represents the motion in 
the mode where both masses are moving in the same direction, at different amplitudes. The matrix for $v_2$ represents 
the case where the masses are moving in opposite directions (indicated by the negative sign), at their respective amplitudes. 

\section{Conclusion}
In this project, we discussed the coupled-mass system, matrix operations, and eigenvalues and eigenvectors. We also worked 
through an example problem with two connected masses in detail to reinforce understanding of the concepts. This foundation 
will aid in our overall goal of analyzing the elasticity of metals.

\section{Bibliography}
\noindent\hangindent=1.5em \hangafter=1
Nagle, R. Kent, Edward B. Saff, and Arthur David Snider (2018). \textit{Fundamentals of Differential Equations}. 9th ed. Pearson. \\
\hspace*{1.5em}\textsc{isbn}: 9780321977113

\noindent\hangindent=1.5em \hangafter=1
Lay, David C., Steven R. Lay, and Judi J. McDonald (2021). \textit{Linear Algebra and Its Applications}. 6th ed. Pearson.

\end{document}